\section{Operating Regime and Composition}
\label{sec:scope}

\policy{} applies to batched positive EOT problems with compatible batched
representations and a two-marginal verifier. Different numbers of iterations
within a batch create removable work; net acceleration depends on runtime as a
function of batch width and on incremental control and state-movement costs
(\cref{eq:active-gate}).

Initialization shapes both the total iteration count and its variation within
a window. The study that uses one initialization for multiple problems uses
controlled overlapping ImageNet-32
supports: dynamic-compaction execution reduces slots by 30.62\% and accelerates the
steady-state critical path by $1.421\times$, including proposal and transfer
costs. The feature-training study uses cold proposals for independently
resampled OT-FM couplings.

The executor schedules problems from completion under the configured stopping
rule. The grouping intervention measures the effect of grouping a fixed problem
multiset using offline oracle iteration numbers. The width-cost sweep separately
measures how the backend converts each survivor width into update time. Together
they instantiate the two terms in \cref{eq:time-difference}; predictive grouping
can use the same quantities when estimates of completion depth are available.

Dynamic compaction trades computation for state movement and buffers.
The matched logical-mask/dynamic-compaction study reports 26.1 MiB of additional peak
allocation alongside lower time and endpoint energy; \cref{tab:phase-costs}
reports its phase costs. OTT-JAX and PyKeOps
realize release through backend-specific rebucketing and batched
independent-problem execution.

The calibrated width-cost curve belongs to a specified device, kernel path,
problem shape, and numerical configuration. The packed Triton path also
switches narrow tails to its matched single-problem solver, so its complete
cost curve includes that routing boundary. The accounting identities and
\cref{eq:time-difference} apply to any measured cost curve; the observed A100
wave locations apply to the reported Triton configuration.
